\BourbakiExercisesSection

\begin{enumerate}
\def\labelenumi{\arabic{enumi}.}
\item
  若 M 是单生成 A-模，则 \(\mathsf{S}(\mathsf{M}) = \mathsf{T}(\mathsf{M})\) 。由此推出一个单同态 \(j: N \to M\) 的例子，使得 \(\mathsf{S}(j): \mathsf{S}(\mathsf{N}) \to \mathsf{S}(\mathsf{M})\) 不是单射（ \(\S 5\) ，习题 1）。
\item
  为对称代数建立与§5习题8(a)和(b)中性质相类似的性质。
\item
  设 \(\mathbf{K}\) 是一个域， \(\mathbf{B}\) 为多项式环 \(\mathbf{K}[\mathbf{X}_1, \mathbf{X}_2]\) ，而 \(\mathfrak{p}\) 为 \(\mathbf{B}\) 中由多项式 \(\mathbf{X}_1^3 - \mathbf{X}_2^2\) 生成的理想。
\end{enumerate}

\begin{enumerate}
\def\labelenumi{(\alph{enumi})}
\item
  证明 p 是素理想，因此 A = B/p 是整环。设 a 和 b 分别表示 \(X_{1}\) 和 \(X_{2}\) 在 A 中的典范像；a 和 b 在 A 中不可逆，且 \(a^{3} = b^{2}\) 。
\item
  设 m 是 A 的由 a 和 b 生成的（极大）理想。证明对称代数 \(\mathsf{S}(\mathsf{m})\) 不是整环。（注意，m 与商模 \(A^{2}/N\) 等同，其中 N 是 \(A^{2}\) 的由 \((b, -a)\) 生成的子模；得出结论：\(\mathsf{S}(\mathsf{m})\) 同构于商环 \(\mathbf{A}[\mathbf{U}, \mathbf{V}]/(b\mathbf{U} - a\mathbf{V})\) ，并证明在多项式环 \(\mathbf{A}[\mathbf{U}, \mathbf{V}]\) 中主理想 \((b\mathbf{V} - a\mathbf{U})\) 不是素理想，这可通过考虑乘积 \(b^{2}(b\mathbf{V}^{3} - \mathbf{U}^{3})\) 来实现。）
\end{enumerate}

¶ 4. 设 A 是一个交换环，a 是 A 的一个理想。理想 a 的里斯代数是单未定元多项式代数 A{[}T{]} 的由形如 \(c_{0} + c_{1}\mathrm{T} + \cdots + c_{n}\mathrm{T}^{n}\) 的多项式组成的子代数 R(a)，其中 \(c_{0} \in \mathbf{A}\) ，且 \(c_{k} \in \mathfrak{a}^{k}\) （对每个整数 \(k \geqslant 1\)）。

\begin{enumerate}
\def\labelenumi{(\alph{enumi})}
\item
  证明存在且仅存在一个满射 A-代数同态 \(r: \mathsf{S}(\mathfrak{a}) \to \mathsf{R}(\mathfrak{a})\) ，使得 \(r\) 在 \(\mathfrak{a}\) 上的限制是 A-线性映射 \(x \to x\mathbb{T}\) 。
\item
  设环 A 是整环。令 \(a_1, \ldots, a_n\) 是 A 的元素，且 \(a_n \neq 0\) ；考虑 A-代数同态
\end{enumerate}

\[
u: \mathrm{A} [ \mathrm{X} _ {1}, \dots , \mathrm{X} _ {n} ] \rightarrow \mathrm{A} [ \mathrm{T} ]
\]

，使得 \(u(\mathbf{X}_i) = a_i\mathrm{T}\) 。核 \(\mathfrak{r} = \operatorname{Ker}(u)\) 是一个分次理想，其齐次分量 \(\mathfrak{r}_m\) （次数为 \(m\) ）由齐次多项式 \(f\) （次数为 \(m\) 且满足 \(f(a_1, \ldots, a_n) = 0\) ）组成。若 \(\mathfrak{r}' \subset \mathfrak{r}\) 是 \(\mathrm{A}[\mathrm{X}_1, \ldots, \mathrm{X}_n]\) 中由齐次分量 \(\mathfrak{r}_1\) （属于 \(\mathfrak{r}\)）生成的理想，证明 A-模 \(\mathfrak{r}/\mathfrak{r}'\) 是一个挠 A-模。（为验证：对所有 \(f \in \mathfrak{r}_m\) ，存在 A 中的 \(c \neq 0\) 使得 \(cf \in \mathfrak{r}'\) ，对 \(m\) 作归纳论证，写出

\[
f \left(\mathrm{X} _ {1}, \dots , \mathrm{X} _ {n}\right) = \mathrm{X} _ {1} f _ {1} \left(\mathrm{X} _ {1}, \dots , \mathrm{X} _ {n}\right) + \mathrm{X} _ {2} f _ {2} \left(\mathrm{X} _ {2}, \dots , \mathrm{X} _ {n}\right) + \dots + \mathrm{X} _ {n} f _ {n} \left(\mathrm{X} _ {n}\right)
\]

其中 \(f_{j}\) 是次数 \(m - 1\) 的齐次多项式；考虑多项式 \(g(\mathbf{X}_1, \ldots, \mathbf{X}_n) = \mathbf{X}_1 f_1(a_1, \ldots, a_n) + \mathbf{X}_2 f_2(a_2, \ldots, a_n) + \cdots + \mathbf{X}_n f_n(a_n)\) ，观察到 \(g \in \mathfrak{r}_1 \subset \mathfrak{r}'\) ，并作差 \(a_n^{m-1} f - \mathbf{X}_n^{m-1} g\) 。

\begin{enumerate}
\def\labelenumi{(\alph{enumi})}
\setcounter{enumi}{2}
\tightlist
\item
  从 (b) 推断，当 A 是整环时，(a) 中同态 r 的核是 S(a) 的挠模。由此推断，为使 S(a) 是整环，必须且只需 S(a) 是无挠 A-模。特别地，若理想 a 是投射 A-模，则 S(a) 是整环且同构于 R(a)。
\end{enumerate}

\begin{enumerate}
\def\labelenumi{\arabic{enumi}.}
\setcounter{enumi}{4}
\tightlist
\item
  设 E 是域 K 上维数有限为 n 的向量空间。
\end{enumerate}

\begin{enumerate}
\def\labelenumi{(\alph{enumi})}
\tightlist
\item
  证明对每个整数 \(m\) ，对称幂 \(S^m(E)\) 与向量空间 \(S_m'(E)\) （即 \(m\) 阶对称逆变张量的空间）具有相同的维数 \(\binom{m + n - 1}{n - 1} = \binom{n + m - 1}{m}\) 。
\end{enumerate}

*(b) 若 K 的特征为 p \textgreater{} 0，证明 p 阶对称化张量所成的向量空间 \(\mathbb{S}_{p}^{\prime\prime}(\mathbf{E})\) 的维数为 \(\binom{n}{p}\) （注意，若在由 E 的 p 个向量组成的张量积 \(x_{1} \otimes x_{2} \otimes \cdots \otimes x_{p}\) 中，对一对不同的指标 i, j 有 \(x_{i} = x_{j}\) ，则该乘积的对称化为零；利用如下事实：若 \((\mathrm{H}_{i})_{1 \leqslant i \leqslant r}\) 是 \(\{1, 2, \ldots, p\}\) 的一个分成非空集合（至少含两个集合）的划分，则 \(S_{p}\) 中保持每个 \(H_{i}\) 稳定的子群的阶不被 p 整除）。另一方面，在 \(\mathbb{S}^{p}(\mathbf{E})\) 中，p 阶对称张量空间 \(\mathbb{S}_{p}^{\prime}(\mathbf{E})\) 的典范像的维数为 n，且由张量 \(x \otimes x \otimes \cdots \otimes x\) （p 次）的像生成，其中 \(x \in E\) （利用同一注记）；因此从 \(\mathbb{S}_{p}^{\prime}(\mathbf{E})\) 到 \(\mathbb{S}^{p}(\mathbf{E})\) 的典范映射不是双射，尽管这两个向量空间具有相同的维数。*

